% GENERATED FILE. Edit canonical lessons, then run npm run generate:books.
% canonical-source-hash: fnv1a64:977127521fb53794
% canonical-lessons: ES-C642-circunvalacion, ES-C642-arcen, ES-C642-peaje, ES-C642-bache, ES-C642-tramo

\chapter{En ruta}
\label{ch:es-en-ruta}
\hlchaptermodality{642}

\begin{chapteropening}
\textbf{By the end of this chapter:} \emph{I can follow a road trip: the ring road, the hard shoulder, potholes, tolls and each stretch of the way.}

Complete the last lesson of chapter 642: I can follow a road trip: the ring road, the hard shoulder, potholes, tolls and each stretch of the way.
\end{chapteropening}

\section[la circunvalación]{la circunvalación — the ring road, the bypass}

\label{lesson:ES-C642-circunvalacion}



\begin{quote}
Before the new one: say the Spanish for the journey, then the Spanish for to stop.
\end{quote}

\subsection*{You'll want to know: la circunvalación}
\textbf{la circunvalación} — \textquotedblleft{}the ring road, the bypass\textquotedblright{}.

\begin{quote}
Por la circunvalación no pasas por el centro.
\end{quote}

\begin{quote}
\emph{On the ring road you do not go through the centre.}
\end{quote}

In some Spanish cities it is called \emph{la ronda}.

\begin{cousinweb}
\emph{circunvalación} is a learned borrowing built on Latin \emph{circumvallāre} \textquotedblleft{}to surround with a wall\textquotedblright{} (Lewis \& Short), from \emph{circum} \textquotedblleft{}around\textquotedblright{} and \emph{vallum} \textquotedblleft{}rampart\textquotedblright{}.
\end{cousinweb}

\begin{grammarlens}[title={What you've built}]
One more word for telling, explaining and describing at B1.
\end{grammarlens}

\subsection*{Your turn}
\begin{itemize}
\raggedright
  \item \emph{Say it:} \emph{la circunvalación}
  \item \emph{Say it:} \emph{la circunvalación}, once more
  \item \emph{Say it:} a sentence of your own with \emph{la circunvalación}
  \item \emph{Recall:} say the Spanish for the journey, then the Spanish for to stop, then say \emph{la circunvalación} again
\end{itemize}

\begin{tcolorbox}[breakable,colback=teal!4,colframe=teal!35!black,title={Before you move on}]
What does la circunvalación mean? (\textquotedblleft{}The ring road, the bypass\textquotedblright{}.) Say it once more.
\end{tcolorbox}
\section[el arcén]{el arcén — the hard shoulder (of a road)}

\label{lesson:ES-C642-arcen}



\begin{quote}
Before the new one: say the Spanish for to stop, then the Spanish for the ring road.
\end{quote}

\subsection*{You'll want to know: el arcén}
\textbf{el arcén} — \textquotedblleft{}the hard shoulder (of a road)\textquotedblright{}.

\begin{quote}
Paramos en el arcén con las luces encendidas.
\end{quote}

\begin{quote}
\emph{We stopped on the hard shoulder with the lights on.}
\end{quote}

A Spain word; Mexico says \emph{el acotamiento} and Argentina \emph{la banquina}.

\begin{cousinweb}
The origin of \emph{arcén} is not settled.
\end{cousinweb}

\begin{grammarlens}[title={What you've built}]
One more word for telling, explaining and describing at B1.
\end{grammarlens}

\subsection*{Your turn}
\begin{itemize}
\raggedright
  \item \emph{Say it:} \emph{el arcén}
  \item \emph{Say it:} \emph{el arcén}, once more
  \item \emph{Say it:} a sentence of your own with \emph{el arcén}
  \item \emph{Recall:} say the Spanish for to stop, then the Spanish for the ring road, then say \emph{el arcén} again
\end{itemize}

\begin{tcolorbox}[breakable,colback=teal!4,colframe=teal!35!black,title={Before you move on}]
What does el arcén mean? (\textquotedblleft{}The hard shoulder (of a road)\textquotedblright{}.) Say it once more.
\end{tcolorbox}
\section[el peaje]{el peaje — the toll}

\label{lesson:ES-C642-peaje}



\begin{quote}
Before the new one: say the Spanish for the ring road, then the Spanish for the hard shoulder.
\end{quote}

\subsection*{You'll want to know: el peaje}
\textbf{el peaje} — \textquotedblleft{}the toll\textquotedblright{}.

\begin{quote}
En esta autopista hay que pagar peaje.
\end{quote}

\begin{quote}
\emph{You have to pay a toll on this motorway.}
\end{quote}

\emph{Una autopista de peaje} — a toll motorway.

\begin{cousinweb}
\emph{peaje} is usually traced to Old French \emph{péage}, from Medieval Latin \emph{pedāticum} \textquotedblleft{}a toll for passing through\textquotedblright{}, built on Latin \emph{pēs, pedis} \textquotedblleft{}a foot\textquotedblright{} (Lewis \& Short).
\end{cousinweb}

\begin{grammarlens}[title={What you've built}]
One more word for telling, explaining and describing at B1.
\end{grammarlens}

\subsection*{Your turn}
\begin{itemize}
\raggedright
  \item \emph{Say it:} \emph{el peaje}
  \item \emph{Say it:} \emph{el peaje}, once more
  \item \emph{Say it:} a sentence of your own with \emph{el peaje}
  \item \emph{Recall:} say the Spanish for the ring road, then the Spanish for the hard shoulder, then say \emph{el peaje} again
\end{itemize}

\begin{tcolorbox}[breakable,colback=teal!4,colframe=teal!35!black,title={Before you move on}]
What does el peaje mean? (\textquotedblleft{}The toll\textquotedblright{}.) Say it once more.
\end{tcolorbox}
\section[el bache]{el bache — the pothole; the bad patch}

\label{lesson:ES-C642-bache}



\begin{quote}
Before the new one: say the Spanish for the hard shoulder, then the Spanish for the toll.
\end{quote}

\subsection*{You'll want to know: el bache}
\textbf{el bache} — \textquotedblleft{}the pothole; the bad patch\textquotedblright{}.

\begin{quote}
Cuidado con los baches.
\end{quote}

\begin{quote}
\emph{Watch out for the potholes.}
\end{quote}

\emph{Pasar un bache} — to go through a rough patch.

\begin{cousinweb}
The origin of \emph{bache} is not settled; it may imitate the thud of a wheel dropping into a hole.
\end{cousinweb}

\begin{grammarlens}[title={What you've built}]
One more word for telling, explaining and describing at B1.
\end{grammarlens}

\subsection*{Your turn}
\begin{itemize}
\raggedright
  \item \emph{Say it:} \emph{el bache}
  \item \emph{Say it:} \emph{el bache}, once more
  \item \emph{Say it:} a sentence of your own with \emph{el bache}
  \item \emph{Recall:} say the Spanish for the hard shoulder, then the Spanish for the toll, then say \emph{el bache} again
\end{itemize}

\begin{tcolorbox}[breakable,colback=teal!4,colframe=teal!35!black,title={Before you move on}]
What does el bache mean? (\textquotedblleft{}The pothole; the bad patch\textquotedblright{}.) Say it once more.
\end{tcolorbox}
\section[el tramo]{el tramo — the stretch, the section (of a road or route)}

\label{lesson:ES-C642-tramo}



\begin{quote}
Before the new one: say the Spanish for the toll, then the Spanish for the pothole.
\end{quote}

\subsection*{You'll want to know: el tramo}
\textbf{el tramo} — \textquotedblleft{}the stretch, the section (of a road or route)\textquotedblright{}.

\begin{quote}
Este tramo está en obras.
\end{quote}

\begin{quote}
\emph{This stretch has roadworks.}
\end{quote}

\emph{Un tramo de escalera} — a flight of stairs.

\begin{cousinweb}
\emph{tramo} is traced in EtymDB to Latin \emph{trāmes} \textquotedblleft{}a cross-way, a by-path\textquotedblright{} (Lewis \& Short). Some sources link it to \emph{trama} \textquotedblleft{}weft\textquotedblright{} instead, so the origin is not settled.
\end{cousinweb}

\begin{grammarlens}[title={What you've built}]
One more word for telling, explaining and describing at B1.
\end{grammarlens}

\subsection*{Your turn}
\begin{itemize}
\raggedright
  \item \emph{Say it:} \emph{el tramo}
  \item \emph{Say it:} \emph{el tramo}, once more
  \item \emph{Say it:} a sentence of your own with \emph{el tramo}
  \item \emph{Recall:} say the Spanish for the toll, then the Spanish for the pothole, then say \emph{el tramo} again
\end{itemize}

\begin{tcolorbox}[breakable,colback=teal!4,colframe=teal!35!black,title={Before you move on}]
What does el tramo mean? (\textquotedblleft{}The stretch, the section (of a road or route)\textquotedblright{}.) Say it once more.
\end{tcolorbox}
